\documentclass[12pt]{article}

\usepackage{sbc-template}
\usepackage{graphicx,url}
\usepackage{float}
\usepackage[utf8]{inputenc}  
\usepackage[brazil]{babel}

\sloppy

\title{Análise Comparativa de Estruturas Disjoint Set Union: \\ Da Abordagem Ingênua à Otimização de Tarjan}

\author{André X. Lazarini, João P. G. Ribeiro, Lucas F. Vilela e Marina F. S. Cabalzar}

\address{Engenharia de software -- Pontifícia Universidade Católica de Minas Gerais (PUC Minas)\\
  Belo Horizonte -- MG -- Brazil
\email{\{1147036, joao.ribeiro.1485762, lucas.vilela, marina.cabalzar\}@sga.pucminas.br}
}

\begin{document} 

\maketitle

\begin{abstract}
  This paper analyzes the impact of optimization heuristics on the Disjoint Set Union (DSU) data structure for solving the dynamic connectivity problem. Three versions were implemented in Java: Naive, Union by Rank, and Full Tarjan. Empirical experiments measured execution time and operation counts across worst-case, balanced, and random scenarios. The results validate the theoretical transition from $O(n)$ complexity to an almost constant amortized time $O(\alpha(n))$, highlighting the necessity of advanced data structures for large-scale algorithms.
\end{abstract}

\begin{resumo} 
  Este artigo analisa o impacto de heurísticas de otimização na estrutura de dados Disjoint Set Union (DSU) para a resolução do problema de conectividade dinâmica. Foram implementadas três variantes da estrutura em Java: Naive, Union by Rank e Full Tarjan. Experimentos empíricos mediram o tempo de execução e a contagem de operações em cenários de pior caso, árvores balanceadas e acessos aleatórios. Os resultados validam a transição teórica da complexidade $O(n)$ para um tempo amortizado quase constante $O(\alpha(n))$, evidenciando a necessidade de estruturas avançadas para algoritmos em larga escala.
\end{resumo}


\section{Introdução}

O gerenciamento eficiente de subconjuntos disjuntos é fundamental para o desempenho de algoritmos de grafos e processamento de imagens. A estrutura Disjoint Set Union (DSU), também conhecida como Union-Find, busca solucionar o problema da conectividade dinâmica. A eficiência dessas operações impacta diretamente a viabilidade de algoritmos como o de Kruskal para Árvores Geradoras Mínimas.

Implementações básicas da DSU frequentemente sofrem com a degradação de desempenho em sequências longas de operações. Sem estratégias de balanceamento, a estrutura interna pode assumir formas lineares, elevando o custo da operação de busca ao nível $O(n)$. Este trabalho investiga como as técnicas introduzidas por Robert Tarjan, especificamente a união por rank e a compressão de caminho, reduzem a complexidade para um tempo amortizado quase constante.

Para comprovar tal ganho de eficiência, este experimento implementa a DSU em três níveis progressivos de otimização: a versão básica (Naive), a otimização por altura (Union by Rank) e a otimização completa (Full Tarjan). Através de uma suíte de benchmarks projetada para forçar cenários adversariais e aleatórios, o presente artigo realiza uma análise comparativa do tempo de execução e da quantidade de operações lógicas alocadas por cada abordagem.

\section{Fundamentação Teórica}

A análise de complexidade da DSU varia drasticamente conforme as heurísticas aplicadas. Na implementação \textit{Naive}, a ausência de critérios de união resulta em árvores desbalanceadas com altura proporcional ao número de elementos $n$. Consequentemente, o limite superior de tempo para a operação \textit{Find} é $O(n)$, tornando a estrutura ineficiente para grandes volumes de dados.

A introdução da união por rank (\textit{Union by Rank}) limita a altura da árvore ao garantir que o conjunto menor seja sempre anexado à raiz do conjunto maior. Essa restrição estrutural assegura que a altura máxima não ultrapasse $\log n$, reduzindo a complexidade para $O(\log n)$. Entretanto, o desempenho ideal é atingido apenas com a adição da compressão de caminho (\textit{Path Compression}).

A técnica de Tarjan utiliza a compressão de caminho para achatar a estrutura da árvore durante cada execução da operação \textit{Find}. A combinação de ambas as heurísticas resulta em uma complexidade amortizada de $O(\alpha(n))$, onde $\alpha$ é a inversa da função de Ackermann. Na prática, essa função cresce de forma tão lenta que seu valor permanece inferior a 5 para qualquer entrada fisicamente computável.

\section{Metodologia}

\subsection{Implementação e Tecnologias}
O projeto foi desenvolvido em Java, com a arquitetura centralizada na interface \texttt{DisjointSet}, que estabelece o contrato das operações fundamentais. O uso de métodos \textit{default} para a instrumentação (contagem de operações) habilitou o uso de polimorfismo na classe de testes (\texttt{Benchmark}), garantindo que o algoritmo avaliador fosse agnóstico em relação à implementação interna da estrutura, facilitando a extensibilidade e blindando os testes contra vieses. O ambiente de execução consistiu em um processador Intel(R) Core(TM) i7-10700F CPU @ 2.90GHz, com 24 GB de RAM, operando sobre o sistema Windows 11. Os testes foram compilados e executados utilizando a JVM OpenJDK 24.0.1 (Temurin).

\subsection{Cenários de Teste}
A avaliação empírica utilizou a classe \texttt{Benchmark} para simular três cenários distintos:
\begin{itemize}
    \item \textbf{Pior Caso (Cadeia Linear):} Força a criação de uma lista encadeada através de uniões sequenciais de $i$ com $i+1$, maximizando a profundidade da busca.
    \item \textbf{Árvore Balanceada:} Constrói uma estrutura perfeitamente equilibrada para observar o comportamento logarítmico.
    \item \textbf{Aleatório:} Simula o uso típico da estrutura com uniões e buscas baseadas em sementes aleatórias fixas para reprodutibilidade.
\end{itemize}

Cada teste foi precedido por uma fase de aquecimento da JVM para estabilizar o compilador JIT e executado em 3 rodadas (\textit{trials}) para mitigar ruídos térmicos ou do sistema operacional.

\section{Discussão de Resultados}

\subsection{Desempenho no Pior Caso (Adversarial)}

O cenário adversarial demonstra o impacto severo da falta de balanceamento. O tempo de execução da versão Naive apresenta um crescimento linear acentuado, alcançando a marca de quase 100 ms para entradas maiores. Em contrapartida, as implementações Union by Rank e Full Tarjan mantêm o tempo de processamento praticamente constante e próximo a 0 ms, evidenciando a eficácia imediata das heurísticas.

\begin{figure}[H]
    \centering
    \includegraphics[width=0.8\textwidth]{1_adversarial_tempo.pdf}
    \caption{Tempo de execução no cenário adversarial (cadeia linear).}
    \label{fig:adv_tempo}
\end{figure}

A análise das operações em memória corrobora o impacto temporal. A contagem de operações por busca evidencia a classe de complexidade $O(n)$ da versão Naive. As abordagens otimizadas mitigam a formação das cadeias longas, achatando a curva de forma estrutural e mantendo os acessos mínimos.

\begin{figure}[H]
    \centering
    \includegraphics[width=0.8\textwidth]{2_adversarial_ops.pdf}
    \caption{Média de operações por busca no cenário adversarial.}
    \label{fig:adv_ops}
\end{figure}

\subsection{Desempenho em Árvores Balanceadas}

No cenário de árvore construída de forma balanceada, a versão Union by Rank exibe um crescimento visivelmente logarítmico, operando paralela à sua complexidade teórica de $O(\log n)$. A adoção da compressão de caminhos na implementação Full Tarjan elimina essa dependência de profundidade logarítmica, sustentando a contagem média de operações em um platô inferior, validando a complexidade $O(\alpha(n))$.

\begin{figure}[H]
    \centering
    \includegraphics[width=0.8\textwidth]{3_balanceada_ops.pdf}
    \caption{Crescimento de operações por busca no cenário balanceado.}
    \label{fig:bal_ops}
\end{figure}

\subsection{Desempenho no Cenário Aleatório}

A simulação de requisições aleatórias reflete cenários de aplicação convencionais de algoritmos de conectividade dinâmica. A versão Naive sofre rápida degradação inicial de tempo e torna-se inviável em larga escala. As versões otimizadas absorvem bem a carga de até 1.000.000 de elementos, contudo, a implementação Full Tarjan consegue entregar tempos médios de resposta ligeiramente inferiores aos da Union by Rank devido ao processo contínuo de encurtamento das árvores em tempo de execução.

\begin{figure}[H]
    \centering
    \includegraphics[width=0.8\textwidth]{4_aleatorio_tempo.pdf}
    \caption{Tempo médio de resposta em operações aleatórias.}
    \label{fig:aleat_tempo}
\end{figure}

\section{Conclusão}

Os experimentos realizados confirmam que a escolha adequada de estruturas de dados é determinante para a escalabilidade de sistemas de software. A implementação básica da DSU mostrou-se inviável para grandes conjuntos de dados, apresentando degradação severa em cenários lineares. Em contrapartida, as heurísticas de Tarjan demonstraram eficiência excepcional, mantendo o desempenho quase constante mesmo com milhões de elementos. A transição da complexidade de $O(n)$ para $O(\alpha(n))$ representa não apenas um ganho quantitativo, mas uma mudança de viabilidade algorítmica essencial para a engenharia de software moderna.

\end{document}